\documentclass[10pt,a4paper]{article}
\usepackage[utf8]{inputenc}
\usepackage{graphicx}
\usepackage[left=1.0in, right=1.0in, top=1.0in, bottom=1.0in]{geometry}
\usepackage[swedish]{babel}
\usepackage{tikz}
\usetikzlibrary{arrows.meta,shapes,calc,backgrounds}
\usepackage{parskip}
\usepackage{xcolor}
\usepackage[most]{tcolorbox}
\usepackage{mdframed}
\usepackage{tabularx}
\usepackage{colortbl}
\usepackage{enumerate}
\usepackage{enumitem}

\setlength{\extrarowheight}{8pt}
\renewcommand{\arraystretch}{1}

% My colors
\definecolor{mycolor1}{rgb}{1,0,0} %red
\definecolor{lightpurple}{HTML}{c69cd6}
\definecolor{darkpurple}{HTML}{7b5c87}

%circled
\newcommand*\circled[1]{\tikz[baseline=(char.base)]{
		\node[shape=circle,draw,minimum size=1em, inner sep=2pt] (char)
		{#1};}
}

% Purpleleftline
\newmdenv[innerleftmargin=20,topline=false,rightline=false,bottomline=false, linewidth=3pt, linecolor=darkpurple]{purpleleftline}

% Kommentarboxen
\tikzstyle{mybox} = [draw=lightpurple!35, fill=lightpurple!35, very thick,
rectangle, inner sep=15pt, inner ysep=15pt]
\tikzstyle{fancytitle} =[fill=darkpurple, text=white]

% Title format
\newcommand{\pagetitle}[2] {
	\begin{tabularx}{\linewidth}{>{\columncolor{darkpurple}}c >{\columncolor{lightpurple!35}}X}
		\LARGE \textcolor{white}{{#1}} & \LARGE {#2} \\[3pt]
	\end{tabularx}
}

% Instruction environment
\newenvironment{instructionsenv}
{
	\vspace{-4pt}
	\begin{purpleleftline}
		\vspace{2em}
		{\Large\bfseries{Instruktioner}}\par
		\vspace{1em}
	}
	{ 
	\end{purpleleftline}
}

% Instruction
\newcommand{\instruction}[2] {
	\textcolor{mycolor1}{\circled{\large {#1}}} 
	{\large #2}\par
}

% Kommentarbox
\newcommand{\kommentarbox}[1]{
	\begin{center}
		\begin{tikzpicture}
		\node [mybox] (box){%
			\begin{minipage}{0.91\textwidth}
			{#1}
			\end{minipage}
		};
		\node[fancytitle, right=15pt] at (box.north west) {\bfseries Kommentar};
		\end{tikzpicture}
	\end{center}	
}

%\begin{bild}{fil}{storlek andel av linewidth}
%\end{bild}
\newenvironment{bild}[2]
{
\begin{center}
	\begin{tikzpicture}
	\node[anchor=south west,inner sep=0] (image) at (0,0) {\includegraphics[width=#2\textwidth]{./img/#1}};
	\begin{scope}[x={(image.south east)},y={(image.north west)}]
}
{		
	\end{scope}
	\end{tikzpicture}
\end{center}
}
%\siffrastreckcirkel{siffra}{x}{y}{storlek cirkel}{x}{y}
\newcommand{\siffrastreckcirkel}[6]{
	\node[mycolor1, thick] (A) at  ({#2},{#3}) {\Large {#1}};
	\node[circle, draw, mycolor1, thick, minimum size=#4cm] (B) at  ({#5},{#6}) {};
	\draw[mycolor1, thick] (A) -- (B);
}

% \siffrastreckellips{siffra}{x}{y}{längd ellips}{höjd ellips}{x}{y}
\newcommand{\siffrastreckellips}[7]{
	\node[mycolor1, thick] (A) at  ({#2},{#3}) {\LARGE {#1}};
	\node[ellipse, draw, mycolor1,thick,minimum height=#5cm,minimum width=#4cm,draw] (B) at  ({#6},{#7}) {};
\draw[mycolor1, thick] (A) -- (B);
}

% \siffrapil{siffra}{från x}{från y}{till x}{till y}
\newcommand{\siffrapil}[5]{
	\node[mycolor1, thick] (A) at  (#2,#3) {\LARGE #1};
	\node (B) at  (#4,#5) {};
\draw[-Stealth, mycolor1, line width=1pt] (A)--(B);
}

\begin{document}

%%%%%%%%%%%%%%%%%%%%%%%%%%%%%%%%%%%%%%%%%%%%%%%%%%%%%%%%%%%%%%
%%% Titelsidan
%%%%%%%%%%%%%%%%%%%%%%%%%%%%%%%%%%%%%%%%%%%%%%%%%%%%%%%%%%%%%%
\pagenumbering{gobble}
\newgeometry{left=0cm,bottom=0cm,right=0cm,top=0cm}
\begin{tikzpicture}

\fill[lightpurple!35] (0,0) rectangle (\paperwidth,\paperheight);

\fill[white] (0,0) rectangle (0.65\paperwidth,\paperheight);

\fill[darkpurple] (0,0.45\paperheight) rectangle (\paperwidth,0.65\paperheight);

\node[align=left,white] at (0.5\paperwidth,0.57\paperheight) 
{\bfseries\Huge Exempel på användande av \\ \bfseries\Huge EMG-utrustningen på HSC};

\node[align=center,white] at (0.5\paperwidth,0.505\paperheight) 
{ \huge En praktisk guide};

\node[align=center,black] at (0.825\paperwidth,0.1\paperheight) 
{\Large Peter Kvillegård \\[3em] Version 1.0 \\ 2020-01-23};

\end{tikzpicture}
\restoregeometry
\newpage

%%%%%%%%%%%%%%%%%%%%%%%%%%%%%%%%%%%%%%%%%%%%%%%%%%%%%%%%%%%%%%
%%% Bakgrund
%%%%%%%%%%%%%%%%%%%%%%%%%%%%%%%%%%%%%%%%%%%%%%%%%%%%%%%%%%%%%%

\pagetitle{Om denna guide}{}\par
\vspace{1em}
Jag och Jonatan Malmström använde under hösten 2019 EMG-utrustningen på HSC i Lund för ett experiment. Vi lärde oss att använda program och apparater av andra elever och av vår handledare. I denna användarguide har vi skrivit ner hur vi använde utrustningen, i förhoppning om att det kan vara användbart för andra.

Utrustningen som användes och beskrivs här är MegaWin 3.01b mjukvara med ???? hårdvara.  

Vårt arbete fokuserade på att få fram medelamplitud under en definierad tid i olika muskler samtidigt, och användarguiden är därför inriktad för det ändamålet. Detta är inte är en allomfattande manual, utan ett exempel på hur utrustningen kan användas. Det är mycket möjligt att det finns bättre och smidigare sätt, och det finns många andra ändamål som inte kommer att beskrivas här.

All text och bilder är tillgängliga under Creative Commons Zero (CC0), vilket lite förenklat innebär att vem som helst har rätt att göra vad som helst med innehållet. Du får använda, modifiera, och distribuera materialet utan begränsningar och utan krav på att attribuera mig. Källkoden är skriven i \LaTeX{} och finns att hitta tillsammans med bilderna på \verb|http://www.github.com/mouboo/lund_emg_userguide|. Om du upptäcker fel eller vill bygga vidare på denna guide är du mycket välkommen att göra så, kontakta mig på \verb|peterkvillegård@gmail.com| för mer information om detta.

\newpage

%%%%%%%%%%%%%%%%%%%%%%%%%%%%%%%%%%%%%%%%%%%%%%%%%%%%%%%%%%%%%%
%%% Innehåll
%%%%%%%%%%%%%%%%%%%%%%%%%%%%%%%%%%%%%%%%%%%%%%%%%%%%%%%%%%%%%%

\pagetitle{Innehåll}{}\par
\vspace{3em}
{\Large

\begin{enumerate}[label={\Roman*},align=left]
\item Logga in \dotfill 1
\item Skapa protokoll \dotfill 2
\item Skapa person \dotfill 3
\item Anslut biomonitor \dotfill 4
\item Kontrollera signalen \dotfill 5
\item Utför mätning \dotfill 6
\item Bearbeta signalen \dotfill 7
\end{enumerate}
}

\newpage
\pagenumbering{arabic}

%%%%%%%%%%%%%%%%%%%%%%%%%%%%%%%%%%%%%%%%%%%%%%%%%%%%%%%%%%%%%%
%%% Logga in | Steg 1/1
%%%%%%%%%%%%%%%%%%%%%%%%%%%%%%%%%%%%%%%%%%%%%%%%%%%%%%%%%%%%%%

%%%%% Title
\pagetitle{I. Logga in}{Steg 1/1}

%%%%% Instruktioner
\begin{instructionsenv}
	\instruction{1}{Kryssa i ``Workstation only''}
	\instruction{2}{Fyll i Username: \texttt{emg}, Password: \texttt{emgemg}}
	\instruction{3}{Klicka på ``OK''}
\end{instructionsenv}

\vspace{50pt}

%%%%% Bilden
\begin{bild}{login}{0.5}
	%\siffrastreckcirkel{siffra}{x}{y}{storlek cirkel}{x}{y}
	\siffrastreckcirkel{1}{0.2}{0.1}{0.5}{0.03}{0.22}	

	%\siffrastreckcirkel{siffra}{x}{y}{storlek cirkel}{x}{y}
	\siffrastreckcirkel{2}{0.5}{0.6}{1.2}{0.22}{0.43}

	% \siffrastreckellips{siffra}{x}{y}{längd ellips}{höjd ellips}{x}{y}
\siffrastreckellips{3}{0.6}{0.25}{1.5}{0.5}{0.39}{0.1}
\end{bild}

\vspace{35pt}

\newpage

%%%%%%%%%%%%%%%%%%%%%%%%%%%%%%%%%%%%%%%%%%%%%%%%%%%%%%%%%%%%%%
%%% II. Skapa protokoll | Steg 1/?
%%%%%%%%%%%%%%%%%%%%%%%%%%%%%%%%%%%%%%%%%%%%%%%%%%%%%%%%%%%%%%

%%%%% Title
\pagetitle{II. Skapa protokoll}{Steg 1/?}

%%%%% Instruktioner
\begin{instructionsenv}
	\instruction{1}{Klicka på ``Protocol''}
\end{instructionsenv}

\vspace{50pt}

%%%%% Bilden
\begin{bild}{3}{0.91}
	%\siffrastreckcirkel{siffra}{x}{y}{storlek cirkel}{x}{y}
	\siffrastreckcirkel{1}{0.25}{0.65}{2}{0.32}{0.85}
\end{bild}

\vspace{35pt}

%%%%% Kommentarboxen
\kommentarbox{
	Du behöver bara skapa protokoll en gång, detta kan sedan användas för alla kommande mätningar.
}
\newpage

%%%%%%%%%%%%%%%%%%%%%%%%%%%%%%%%%%%%%%%%%%%%%%%%%%%%%%%%%%%%%%
%%% II. Skapa protokoll | Steg 2/?
%%%%%%%%%%%%%%%%%%%%%%%%%%%%%%%%%%%%%%%%%%%%%%%%%%%%%%%%%%%%%%

%%%%% Title
\pagetitle{II. Skapa protokoll}{Steg 2/?}

%%%%% Instruktioner
\begin{instructionsenv}
	\instruction{1}{Klicka på ``New...''}
\end{instructionsenv}

\vspace{50pt}

%%%%% Bilden
\begin{bild}{7}{0.8}
	% \siffrastreckellips{siffra}{x}{y}{längd ellips}{höjd ellips}{x}{y}
	\siffrastreckellips{1}{0.75}{0.75}{1.5}{0.6}{0.92}{0.695}
\end{bild}

\vspace{35pt}
\newpage

%%%%%%%%%%%%%%%%%%%%%%%%%%%%%%%%%%%%%%%%%%%%%%%%%%%%%%%%%%%%%%
%%% II. Skapa protokoll | Steg 3/?
%%%%%%%%%%%%%%%%%%%%%%%%%%%%%%%%%%%%%%%%%%%%%%%%%%%%%%%%%%%%%%

%%%%% Title
\pagetitle{II. Skapa protokoll}{Steg 3/?}

%%%%% Instruktioner
\begin{instructionsenv}
	\instruction{1}{Fyll i protokollets namn}
	\instruction{2}{Välj ``Advanced protocol''}
	\instruction{3}{Klicka på ``Next''}
\end{instructionsenv}

\vspace{50pt}

%%%%% Bilden
\begin{bild}{8}{0.8}
	%\siffrastreckcirkel{siffra}{x}{y}{storlek cirkel}{x}{y}	
	\siffrastreckcirkel{1}{0.55}{0.7}{1}{0.25}{0.8}

	% \siffrastreckellips{siffra}{x}{y}{längd ellips}{höjd ellips}{x}{y}
	\siffrastreckellips{2}{0.5}{0.55}{3.5}{0.8}{0.15}{0.565}

	% \siffrastreckellips{siffra}{x}{y}{längd ellips}{höjd ellips}{x}{y}
\siffrastreckellips{3}{0.5}{0.15}{2.5}{0.7}{0.495}{0.037}
\end{bild}

\vspace{35pt}
%%%%% Kommentarboxen
\kommentarbox{
	I denna guide har jag valt att kalla protokollet ``ledare'', döp gärna ert protokoll till något unikt som ni kommer ihåg, exempelvis försöksledarnas initialer.
}
\newpage

%%%%%%%%%%%%%%%%%%%%%%%%%%%%%%%%%%%%%%%%%%%%%%%%%%%%%%%%%%%%%%
%%% II. Skapa protokoll | Steg 4/?
%%%%%%%%%%%%%%%%%%%%%%%%%%%%%%%%%%%%%%%%%%%%%%%%%%%%%%%%%%%%%%

%%%%% Title
\pagetitle{II. Skapa protokoll}{Steg 4/?}

%%%%% Instruktioner
\begin{instructionsenv}
	\instruction{1}{Välj ``ME6000-T16''}
	\instruction{2}{Klicka på ``Next''}
\end{instructionsenv}

\vspace{50pt}

%%%%% Bilden
\begin{bild}{9}{0.8}
	% \siffrastreckellips{siffra}{x}{y}{längd ellips}{höjd ellips}{x}{y}
	\siffrastreckellips{1}{0.55}{0.8}{2}{0.4}{0.2}{0.805}

	% \siffrastreckellips{siffra}{x}{y}{längd ellips}{höjd ellips}{x}{y}
	\siffrastreckellips{2}{0.5}{0.15}{2.5}{0.7}{0.495}{0.037}
\end{bild}

\vspace{35pt}

\newpage

%%%%%%%%%%%%%%%%%%%%%%%%%%%%%%%%%%%%%%%%%%%%%%%%%%%%%%%%%%%%%%
%%% II. Skapa protokoll | Steg 5/?
%%%%%%%%%%%%%%%%%%%%%%%%%%%%%%%%%%%%%%%%%%%%%%%%%%%%%%%%%%%%%%

%%%%% Title
\pagetitle{II. Skapa protokoll}{Steg 5/?}

%%%%% Instruktioner
\begin{instructionsenv}
	\instruction{1}{Välj ``General body'', eller annat alternativ}
	\instruction{2}{Klicka på ``Next''}
\end{instructionsenv}

\vspace{50pt}

%%%%% Bilden
\begin{bild}{10}{0.8}
	% \siffrastreckellips{siffra}{x}{y}{längd ellips}{höjd ellips}{x}{y}
	\siffrastreckellips{1}{0.5}{0.72}{4.5}{1.5}{0.22}{0.73}

	% \siffrastreckellips{siffra}{x}{y}{längd ellips}{höjd ellips}{x}{y}
	\siffrastreckellips{2}{0.5}{0.15}{2.5}{0.7}{0.495}{0.037}
\end{bild}

\vspace{35pt}

%%%%% Kommentarboxen
\kommentarbox{
	Du kan även välja att begränsa vilka muskler som ska visas genom att välja ett annat alternativ. Vi hittade inget sätt att skapa ett protokoll som använder muskler från olika kategorier här, därför rekommenderar vi att välja ``General body'', eller den kategorin som passar bäst för ert experiment, och vid behov lägga till muskler som saknas i kommande steg.
}

\newpage

%%%%%%%%%%%%%%%%%%%%%%%%%%%%%%%%%%%%%%%%%%%%%%%%%%%%%%%%%%%%%%
%%% II. Skapa protokoll | Steg 6/?
%%%%%%%%%%%%%%%%%%%%%%%%%%%%%%%%%%%%%%%%%%%%%%%%%%%%%%%%%%%%%%

%%%%% Title
\pagetitle{II. Skapa protokoll}{Steg 6/?}

%%%%% Instruktioner
\begin{instructionsenv}
	\instruction{1}{Markera första kanalen}
	\instruction{2}{Klicka på ``Select source...''}
\end{instructionsenv}

\vspace{50pt}

%%%%% Bilden
\begin{bild}{11}{0.8}
	% \siffrapil{siffra}{från x}{från y}{till x}{till y}
	\siffrapil{1}{0.2}{0.79}{0.6}{0.79}

	% \siffrastreckellips{siffra}{x}{y}{längd ellips}{höjd ellips}{x}{y}
	\siffrastreckellips{2}{0.2}{0.67}{3}{1}{0.2}{0.53}
\end{bild}

\vspace{35pt}

%%%%% Kommentarboxen
\kommentarbox{
	Kanalerna motsvarar samma ordning som sladdarna är märkta. ME6000-T16 har 4 ingångar som vardera kan ha 4 kanaler. 
}

\newpage

%%%%%%%%%%%%%%%%%%%%%%%%%%%%%%%%%%%%%%%%%%%%%%%%%%%%%%%%%%%%%%
%%% II. Skapa protokoll | Steg 7/?
%%%%%%%%%%%%%%%%%%%%%%%%%%%%%%%%%%%%%%%%%%%%%%%%%%%%%%%%%%%%%%

%%%%% Title
\pagetitle{II. Skapa protokoll}{Steg 7/?}

%%%%% Instruktioner
\begin{instructionsenv}
	\instruction{1}{Välj muskeln är på höger eller vänster sida}
	\instruction{2}{Markera rätt muskel}
	\instruction{3}{Klicka på ``OK''}
\end{instructionsenv}

\vspace{50pt}

%%%%% Bilden
\begin{bild}{12}{0.8}
	% \siffrapil{siffra}{från x}{från y}{till x}{till y}
	\siffrapil{1}{0.2}{0.79}{0.6}{0.79}

	% \siffrastreckellips{siffra}{x}{y}{längd ellips}{höjd ellips}{x}{y}
	\siffrastreckellips{2}{0.2}{0.67}{3}{1}{0.2}{0.53}
\end{bild}

\vspace{35pt}

%%%%% Kommentarboxen
\kommentarbox{
	Kanalerna motsvarar samma ordning som sladdarna är märkta. ME6000-T16 har 4 ingångar som vardera kan ha 4 kanaler. 
}

\newpage

\end{document}